\section{重要研究成果}\label{sec:results}

\subsection{弱哥德巴赫猜想的解决}

\subsubsection{维诺格拉多夫定理（1937）}

维诺格拉多夫使用圆法证明了三素数定理\cite{vinogradov1937representation}：

\begin{theorem}[三素数定理]
    每个充分大的奇数都可以表示为三个素数之和。
\end{theorem}

该定理的证明基于以下关键估计：

\begin{equation}
    r_3(n) = \frac{n^2}{2(\log n)^3} \mathfrak{S}(n) + O\left(\frac{n^2}{(\log n)^4}\right)
\end{equation}

其中 $\mathfrak{S}(n)$ 是奇异级数：
\begin{equation}
    \mathfrak{S}(n) = \prod_{p|n} \left(1 - \frac{1}{(p-1)^2}\right) \prod_{p \nmid n} \left(1 + \frac{1}{(p-1)^3}\right)
\end{equation}

\subsubsection{黑尔弗戈特的完全证明（2013）}

黑尔弗戈特通过改进圆法并结合分区的精细估计，证明了弱哥德巴赫猜想对\textbf{所有}大于 $5$ 的奇数成立\cite{helfgott2013ternary}（此前维诺格拉多夫定理只覆盖充分大的奇数，剩余有限范围由数值验证补足）。至此弱哥德巴赫猜想得到完全解决。作为定量的补充，陶哲轩证明了每个大于 $1$ 的奇数都是至多五个素数之和\cite{tao2014hardy}。

\subsection{强哥德巴赫猜想的进展}

\subsubsection{布伦定理（1920）}

布伦使用布伦筛法证明了\cite{brun1920sieve}：

\begin{theorem}[布伦定理]
    每个充分大的偶数都可以表示为两个殆素数之和，即两个素因子个数不超过9的数之和。
\end{theorem}

布伦的常数经后人不断改进：Ramaré（1995）证明了每个偶数都是至多 $6$ 个素数之和\cite{ramare1995explicit}，距离陈氏定理的 $P_2$ 形式仅一步之遥。

\subsubsection{陈氏定理（1966）}

陈景润证明了哥德巴赫猜想研究中最强的结果：

\begin{theorem}[陈氏定理]
    每个充分大的偶数都可以表示为一个素数和一个不超过两个素数乘积之和，即 $n = p + P_2$。
\end{theorem}

（陈景润的原始结果发表于1966年\cite{chen1966representation}，1973年给出完整证明。）

陈氏定理的证明综合使用了筛法和解析数论技术。其核心步骤包括：

\begin{enumerate}
    \item 使用加权筛法构造辅助函数
    \item 应用布伦-蒂奇马什定理估计某种类型的整数分布
    \item 通过精细的分析得到最终结果
\end{enumerate}

\subsection{数值验证}

\begin{table}[htbp]
    \centering
    \caption{哥德巴赫猜想的数值验证历史}
    \label{tab:verification}
    \begin{tabular}{@{}lll@{}}
        \toprule
        \textbf{年份} & \textbf{验证范围} & \textbf{研究者} \\
        \midrule
        1938 & $10^5$ & 皮平（Pipping） \\
        1965 & $10^7$ & 斯坦与斯坦（Stein \& Stein） \\
        1998 & $10^{14}$ & 德舒耶与德莱斯（Deshouillers--te Riele--Saouter） \\
        2001 & $4 \times 10^{14}$ & 里奇斯坦（Richstein） \\
        2014 & $4 \times 10^{18}$ & 奥利维拉·e·席尔瓦等（Oliveira e Silva 等） \\
        \bottomrule
    \end{tabular}
\end{table}

截至目前的最新验证记录为 $4 \times 10^{18}$（Oliveira e Silva--Herzog--Pardi，2014）\cite{oliveira2014computational}。

\subsection{相关变体问题}

\subsubsection{哥德巴赫数}

\begin{definition}[哥德巴赫数]
    如果偶数 $n$ 可以表示为两个素数之和，则称 $n$ 为哥德巴赫数。
\end{definition}

\subsubsection{殆哥德巴赫猜想}

殆哥德巴赫猜想研究的是几乎所有的偶数是否都是哥德巴赫数。

\begin{theorem}[殆哥德巴赫猜想]
    设 $E(x)$ 表示不超过 $x$ 的非哥德巴赫偶数个数，则
    \begin{equation}
        E(x) \ll x^{1-\delta}
    \end{equation}
    对某个 $\delta > 0$ 成立。
\end{theorem}
